\BeleskeID{09_2-s057}
% element: 09_2-s057:e05
Smanjenje kašnjenja može se postići smanjenjem ekvivalentne otpornosti ili kapacitivnosti. Izraz za P tranzistor povezuje njegovu ekvivalentnu otpornost sa naponom napajanja i strujom zasićenja.

U prikazanom objašnjenju povećanje napajanja smanjuje dinamičku otpornost, ali povećava disipaciju sa trećim stepenom, pa zahteva znatno bolje hlađenje procesora. Overklokovanje bez povećanja napajanja neće uspeti, a bez boljeg hlađenja može dovesti do stradanja procesora.

Kod baterijskog napajanja važniji je ograničen raspoloživ izvor energije. Kada puna brzina nije potrebna, smanjuju se učestanost i napajanje, čime se štedi energija i produžava rad uređaja. Ova tehnika promene napona i učestanosti koristi se u mobilnim telefonima, laptopovima i sličnim sistemima.
